% Generated draft: requires mathematical and source review.
\documentclass[11pt]{article}
\usepackage[margin=25mm]{geometry}
\usepackage{fontspec}
\setmainfont{Latin Modern Roman}
\usepackage{amsmath,amssymb,mathtools,unicode-math}
\setmathfont{Latin Modern Math}
\usepackage{enumitem,needspace,xurl,hyperref}
\setlength{\parindent}{0pt}
\setlength{\parskip}{5pt}
\setlength{\emergencystretch}{3em}
\newcommand{\N}{\mathbb{N}}
\newcommand{\Z}{\mathbb{Z}}
\newcommand{\Q}{\mathbb{Q}}
\newcommand{\R}{\mathbb{R}}
\newcommand{\C}{\mathbb{C}}
\newcommand{\F}{\mathbb{F}}
\newcommand{\A}{\mathbb{A}}
\newcommand{\PP}{\mathbb{P}}
\newcommand{\E}{\mathbb{E}}
\newcommand{\eps}{\varepsilon}
\newcommand{\dd}{\,\mathrm{d}}
\begin{document}
\section*{UnsolvedMath ID: 700546}
\textbf{Source problem code:} catalogue entry
\subsection{Definitions and mathematical statement}
Let \(G\) be a group, with no finiteness assumption, and let \(A\) be a nontrivial abelian minimal normal subgroup of \(G\). Thus \(A\trianglelefteq G\), \(A\neq\{1\}\), and every subgroup \(B\leq A\) that is normal in \(G\) satisfies \(B=\{1\}\) or \(B=A\).

A group \(H\) is soluble if its derived series, defined by \(H^{(0)}=H\) and \(H^{(i+1)}=[H^{(i)},H^{(i)}]\), reaches \(\{1\}\) after finitely many steps. Interpret finite rank as finite Pr\"ufer rank: there exists an integer \(r\geq 0\) such that every finitely generated subgroup of \(H\) can be generated by at most \(r\) elements. The supplied record does not explicitly define its rank convention; this interpretation requires confirmation against the source.

An elementary abelian \(p\)-group is an abelian group in which every nonidentity element has order \(p\), where \(p\) is prime; equivalently, it is the additive group of a vector space over \(\F_p\), with no restriction on dimension.

\textbf{Question.} For every such pair \((G,A)\), if \(G/A\) is soluble and has finite rank, must there exist a prime \(p\) such that \(A\) is an elementary abelian \(p\)-group?

\textbf{Supplied qualification.} The source reports that the answer is affirmative under the additional assumption that \(G/A\) is locally polycyclic, and attributes this result to D. I. Zaitsev (1980). Here locally polycyclic means that every finitely generated subgroup is polycyclic; a group is polycyclic if it admits a finite subnormal series with cyclic factors. This reported special case is supplied background, not an independently verified result in this draft.
\subsection{Short English statement}
Must every nontrivial abelian minimal normal subgroup of a group be elementary abelian for some prime whenever the quotient is soluble of finite rank? Finite rank is interpreted here as finite Pr\"ufer rank, subject to source confirmation. The supplied source reports an affirmative answer when the quotient is additionally locally polycyclic.
\subsection{Sources}
\begin{itemize}
\item Local UnsolvedMath registry, permanent ID \(700546\), exact source code \texttt{catalogue entry}, title \emph{Kourovka Notebook Problem 9.11}. The supplied registry labels the problem open; that label has not been independently verified.
\item \emph{The Kourovka Notebook (Unsolved Problems in Group Theory)}, original Problem 9.11, ninth issue (1984); supplied edition reference: twenty-first issue edition (2026), arXiv:1401.0300v48. Proposed by D. I. Zaitsev. The supplied collection description names E. I. Khukhro and V. D. Mazurov as editors.
\item D. I. Zaitsev, \emph{Algebra and Logic}, volume 19, no. 2 (1980), pp. 94--105. Supplied reference for the affirmative answer under the additional locally polycyclic quotient hypothesis. No article title is supplied, and the reference and reported result have not been independently checked.
\end{itemize}
\end{document}
